% ============================================================================
%  Chapter 1 -- PRG32 and PRG32-QT: Setting Up Your Workshop     (Lecture 01)
% ============================================================================
\chapter{PRG32 and PRG32-QT: Setting Up Your Workshop}
\label{ch:setup}

\epigraphline{Give me six hours to chop down a tree and I will spend the first
four sharpening the axe.}{attributed to Abraham Lincoln}

\begin{chapgoals}
By the end of this chapter you will have
\begin{itemize}
\item a clear picture of the pieces of \prg: runtime, cartridge, tools, hosts, Store;
\item \prgqt installed on your laptop, to play any cartridge and look inside it;
\item the \riscv toolchain installed and checked, to build cartridges of your own;
\item one cartridge built by you, running in \prgqt and stopped in the debugger;
\item a list of what to do when something does not work.
\end{itemize}
\end{chapgoals}

A craftsman's first job in a new workshop is to put the tools where the hand
expects them. This chapter is that job. It contains no theory, and it is the
only chapter of the book whose content will age: version numbers are those of
October 2026, and the release pages of the projects should be checked before
following the steps. Everything else in the book depends on the last line of
the summary above: the loop of building a cartridge, sending it to a host and
stopping it in the debugger is repeated in every chapter that follows.

The commands of Section~\ref{sec:setup:build} were run for this book on macOS
against \prgqt. The installations for Windows and Linux follow the official
instructions of the projects and were not run for this text.

% ----------------------------------------------------------------------------
\section{Four levels: install what you need}

Not everything has to be installed, and not everything at once.
Table~\ref{tab:setup:levels} orders the pieces by need. The course, and the
book, can be followed with levels 1 and 2 only: \emph{no board is required}.

\begin{table}[htb]
\centering\small
\caption{Four levels of installation.}\label{tab:setup:levels}
\begin{tabularx}{\textwidth}{@{}cP{0.24\textwidth}LP{0.22\textwidth}@{}}
\toprule
\textbf{Level} & \textbf{You install} & \textbf{You can} & \textbf{Needed}\\
\midrule
1 & \prgqt & play cartridges, download from the Store, use the debugger & by everybody, at once\\
2 & ESP-IDF and the \prg repository & build your own cartridges from assembly or C & by everybody, before Chapter~\ref{ch:firstcontact}\\
3 & the QEMU emulator & run the real firmware on your laptop & optional\\
4 & the kit: an ESP32-C6 board & run on real hardware & optional\\
\bottomrule
\end{tabularx}
\end{table}

Level~1 takes five minutes. Level~2 is a large download and takes from half
an hour to an hour: it is best done at home, on a good connection. Levels 3
and 4 are for the curious, and for the laboratory.

% ----------------------------------------------------------------------------
\section{The platform, in detail}

\subsection{A runtime, a file format, a set of tools}

The name \prg covers four things.

\begin{itemize}
\item A \textbf{resident firmware}, also called the \term{runtime}: the program
  that owns the screen, the buttons and the sound.
\item A \textbf{cartridge format}: a game as one small file, \fname{.prg32}.
\item A \textbf{tool} on the laptop: the command \code{python3 -m prg32}, which
  builds, uploads and packages.
\item A \textbf{Cartridge Store}, where finished games are shared.
\end{itemize}

The code of a cartridge runs as real 32-bit \riscv machine code, RV32IMAC, and
the contract between a game and the runtime is three functions: \code{init},
\code{update} and \code{draw}. The runtime calls \code{init} once and then
\code{update} and \code{draw} at every frame; Chapter~\ref{ch:firstcontact}
starts from this contract.

\subsection{From a source file to the screen}

Figure~\ref{fig:setup:flow} follows a program from the editor to the screen.
The source, in assembly or in C, is translated into machine code by the
\riscv toolchain; the compiler is \code{riscv32-esp-elf-gcc}, installed by
ESP-IDF, the development kit of Espressif. The result is packed into a
cartridge, made of code, data and a header. The cartridge is sent over the
network to a \term{host}, where the runtime loads it and begins to call its
three functions.

\begin{figure}[htb]
\centering
\begin{tikzpicture}[node distance=5mm]
  \node[blkacc,text width=21mm,minimum height=15mm] (a) {\texttt{verm1.S}\\[1pt]\footnotesize your source};
  \node[blkamb,text width=21mm,minimum height=15mm,right=of a] (b) {compiler\\[1pt]\footnotesize \riscv toolchain};
  \node[blkvio,text width=23mm,minimum height=15mm,right=of b] (c) {\texttt{verm1.prg32}\\[1pt]\footnotesize a cartridge};
  \node[blkgrn,text width=19mm,minimum height=15mm,right=of c] (d) {upload\\[1pt]\footnotesize to a host};
  \node[blkred,text width=21mm,minimum height=15mm,right=of d] (e) {runtime\\[1pt]\footnotesize calls \texttt{init}, \texttt{update}, \texttt{draw}};
  \foreach \x/\y in {a/b,b/c,c/d,d/e} \draw[->,thick] (\x) -- (\y);
  \draw[decorate,decoration={brace,mirror,amplitude=4pt}] ([yshift=-2mm]a.south west) -- node[below=5pt,lbl]{on your laptop: needs level 2} ([yshift=-2mm]c.south east);
  \draw[decorate,decoration={brace,mirror,amplitude=4pt}] ([yshift=-2mm]d.south west) -- node[below=5pt,lbl]{a host: \prgqt, QEMU or the board} ([yshift=-2mm]e.south east);
\end{tikzpicture}
\caption{From a source file to the screen.}\label{fig:setup:flow}
\end{figure}

A cartridge is loaded at address \hex{40800000} and reaches the runtime only
through a table of functions, the \term{application binary interface} (ABI).
This is the reason why the same file runs on three different hosts: it never
touches the hardware directly, and each host supplies the same table.

\subsection{The kit: a console on a breadboard}

The hardware of \prg (Figure~\ref{fig:setup:kit}) is built around an Espressif
ESP32-C6: one \riscv core at 160~MHz, with Wi-Fi. A 2.8~inch colour display of
320 by 240 pixels (an ILI9341 controller on an SPI bus) shows the games, which
draw in a viewport of 320 by 200. One or two small audio amplifiers
(MAX98357A, on an I2S bus), a joystick and buttons are wired straight to the
pins of the chip. A game reaches the board over Wi-Fi: the board creates a
network called \code{PRG32}, and the cartridge is uploaded to it. The bill of
materials and the wiring are in the folder \fname{docs/hardware} of the
repository.

\begin{figure}[htb]
\centering
\includegraphics[width=0.7\textwidth]{kit_breadboard.png}
\caption{The \prg kit on a breadboard.}\label{fig:setup:kit}
\end{figure}

\subsection{Three ways to run the same cartridge}

The cartridge file is the same on every host; what changes is how its
instructions are executed and what can be observed
(Table~\ref{tab:setup:hosts}).

\begin{table}[htb]
\centering\small
\caption{Three hosts for the same cartridge.}\label{tab:setup:hosts}
\begin{tabularx}{\textwidth}{@{}P{0.19\textwidth}LLL@{}}
\toprule
 & \textbf{ESP32-C6 board} & \textbf{Espressif QEMU} & \textbf{\prgqt}\\
\midrule
How code runs & natively on the chip & natively on an emulated ESP32-C3 & interpreted RV32IMAC guest\\
Needs a toolchain & yes (ESP-IDF) & yes (ESP-IDF) & no, to play and debug\\
Cartridge arrives by & Wi-Fi upload, Store & staging into the flash image & file import, Store, HTTP upload\\
Instruction stepping & with JTAG and GDB & with GDB & built-in debugger\\
Runs on & the hardware kit & Windows, Linux, macOS & desktop, Raspberry Pi, iOS, Android, TV\\
\bottomrule
\end{tabularx}
\end{table}

The rule of thumb is: use the board when you want the truth about
\emph{time}, and \prgqt when you want the truth about a \emph{register}. The
emulator models an ESP32-C3 rather than a C6 because Espressif maintains a
mature emulator with graphics for that chip; it executes the same \riscv
instructions, and the cartridge does not notice.

\subsection{PRG32-QT: the console on the laptop}

\prgqt is an application for Windows, macOS, Linux, Raspberry Pi, telephones,
tablets and televisions. It contains an interpreter of RV32IMAC, which
executes a cartridge instruction by instruction and offers it the same
functions as the firmware. It has three speeds, chosen in \emph{Settings}:
\emph{Accurate}, timed like the ESP32-C6, which is the default;
\emph{Optimal}, at 30 frames per second; and \emph{Unlimited}, as fast as the
laptop goes. The accurate profile charges each instruction and each call a
cost in cycles of a 160~MHz clock: it is not a cycle-exact model of the chip,
but it is close enough to see when a frame is late. \prgqt also contains a
browser of the Store and the debugger described in Chapter~\ref{ch:intro}.

The controls of the console map onto the keyboard as in
Table~\ref{tab:setup:keys}. A USB or Bluetooth game controller is recognised
by \prgqt when it is plugged in.

\begin{table}[htb]
\centering\small
\caption{The controls, on a keyboard.}\label{tab:setup:keys}
\begin{tabular}{@{}lll@{}}
\toprule
\textbf{On the kit} & \textbf{In \prgqt} & \textbf{In QEMU}\\
\midrule
Joystick & arrow keys, or \key{W A S D} & the same, typed in the terminal\\
Button A & \key{Z} or \key{J} & \key{Z} or \key{J}\\
Button B & \key{X} or \key{K} & \key{X} or \key{K}\\
Start / Select & \key{Enter} or \key{Space} & \key{Enter} or \key{Space}\\
Leave the game & \key{Escape}: back to Setup & close QEMU\\
\bottomrule
\end{tabular}
\end{table}

\begin{prgwarn}
QEMU reads the keyboard from the terminal that launched it, not from the
window with the picture. It is the classic surprise of the first session.
\end{prgwarn}

% ----------------------------------------------------------------------------
\section{Step 1: install PRG32-QT}

Five minutes; no compiler, no account. From the page
\url{https://github.com/riscv-prg32/PRG32-QT}, under \emph{Releases}, take the
latest version of the file for your system (Table~\ref{tab:setup:files}). A
file \fname{SHA256SUMS-platforms.txt} on the same page lists the checksums; builds for
Android and for televisions are there too.

\begin{table}[htb]
\centering\small
\caption{Installers of \prgqt (version 0.4.1, October 2026; the names follow
this pattern with the current number).}\label{tab:setup:files}
\begin{tabular}{@{}lll@{}}
\toprule
\textbf{System} & \textbf{File} & \textbf{Kind}\\
\midrule
Windows 10 or 11 & \fname{PRG32-QT-0.4.1-Windows-X64-Setup.exe} & installer\\
macOS, Apple Silicon & \fname{PRG32-QT-0.4.1-macOS-ARM64.dmg} & disk image\\
macOS, Intel & \fname{PRG32-QT-0.4.1-macOS-Intel-X64.dmg} & disk image\\
Linux, 64-bit PC & \fname{PRG32-QT-0.4.1-Linux-X64.deb} & Debian package\\
Raspberry Pi OS, 64-bit & \fname{PRG32-QT-0.4.1-Raspberry-Pi-OS-ARM64.deb} & Debian package\\
\bottomrule
\end{tabular}
\end{table}

\paragraph{Windows.} Download the file ending in
\fname{Windows-X64-Setup.exe} and run it. If Windows shows a blue SmartScreen
warning, choose \emph{More info} and then \emph{Run anyway}. Accept the
proposed folder: the installer brings the Qt libraries and the compiler
runtime, and nothing else has to be installed. Start \emph{PRG32} from the
Start menu; if the firewall asks, allow it on private networks. The question
appears because \prgqt listens on port 8080 for cartridges sent by the tools.

\paragraph{macOS.} From the Apple menu choose \emph{About This Mac}: ``Chip
Apple M\ldots'' means ARM64, ``Processor Intel'' means X64. Open the
\fname{.dmg} and drag \emph{PRG32} onto the Applications folder. At the first
start, if macOS says that the developer cannot be verified (the application is
distributed outside the App Store), right-click \emph{PRG32} and choose
\emph{Open}, or go to \emph{System Settings, Privacy and Security, Open
Anyway}. Accept incoming connections and local-network access when asked:
they serve the device interface on port 8080 and its announcement on the local
network.

\paragraph{Linux.} The package needs Qt~6.5 or newer, which Debian~13 and
Ubuntu~25.04 provide; on older systems (Ubuntu~24.04 has Qt~6.4) the program
is built from source (Listing~\ref{lst:setup:qtlinux}). On Raspberry Pi OS use
the ARM64 package, or the script \fname{scripts/build-raspbian.sh}.

\shfile[caption={\prgqt on Linux.},label=lst:setup:qtlinux]{code/l01/c06.sh}

\begin{prgtry}[Run it: your first cartridge, from the Store]
\begin{enumerate}[nosep]
\item Start \prgqt: the \emph{Setup} screen appears, the same that the
  firmware shows on the kit.
\item Choose \emph{Browse Store}: the catalogue of the Cartridge Store is
  downloaded.
\item Pick a game, download it, and choose \emph{Run Cartridge}.
\item Play it with the arrow keys, \key{Z} and \key{X}. \key{Escape} brings
  you back to Setup.
\item Open \emph{Settings}: look at the three performance profiles and at the
  layout choices.
\end{enumerate}
\textbf{Expected result.} A game from the Store runs on your laptop
(Figure~\ref{fig:setup:screens}). This is level~1: from now on you can run any
cartridge, and look inside it. If the list of the Store stays empty, the
laptop is not reaching the Internet or a proxy is in the way: \emph{Settings}
shows the address of the Store.
\end{prgtry}

\begin{figure}[htb]
\centering
\begin{subfigure}[b]{0.32\textwidth}\includegraphics[width=\textwidth]{menu_main.png}\caption{The Setup menu}\end{subfigure}\hfill
\begin{subfigure}[b]{0.32\textwidth}\includegraphics[width=\textwidth]{prg32qt_store.png}\caption{Browse Store}\end{subfigure}\hfill
\begin{subfigure}[b]{0.32\textwidth}\includegraphics[width=\textwidth]{prg32qt_settings.png}\caption{Settings}\end{subfigure}
\caption{What you should see: the same three screens on every system.}
\label{fig:setup:screens}
\end{figure}

% ----------------------------------------------------------------------------
\section{Step 2: install the toolchain}

\subsection{One idea behind every installation}

The compiler for the ESP32 family comes with ESP-IDF, the development kit of
Espressif \citep{espidf}. \prg needs ESP-IDF~5.4 or newer, installed for two
chips: \code{esp32c6}, the kit, and \code{esp32c3}, the emulator. After that,
one command does everything: \code{python3 -m prg32}.

\begin{prgwarn}[ESP-IDF is not active by itself]
Every new terminal must switch ESP-IDF on before the tools work. Almost every
installation problem is ESP-IDF not being active \emph{in the terminal you are
typing in}. The symptom is always the same message:
\begin{center}\code{idf.py: command not found}\end{center}
(or \code{riscv32-esp-elf-gcc} not found). On Windows, activating means
opening the \emph{ESP-IDF PowerShell}; on macOS and Linux, running
\fname{export.sh}.
\end{prgwarn}

\subsection{Windows}

Install Git for Windows and Python~3, then run the ESP-IDF Tools Installer and
work in the PowerShell window that it creates
(Listing~\ref{lst:setup:win}). Use the ESP-IDF PowerShell, never a plain
terminal; on Windows the command is \code{python}, not \code{python3};
\code{pip} installs into the environment of ESP-IDF, which is intended. Avoid
folders with spaces or accented letters in their path. Those who prefer an
editor may use Visual Studio Code with the Espressif ESP-IDF extension, which
installs the same toolchain.

\shfile[caption={The toolchain on Windows.},label=lst:setup:win]{code/l01/c09.sh}

\subsection{macOS}

Homebrew comes first, because it provides \code{git}, \code{cmake} and
\code{python} (Listing~\ref{lst:setup:mac}); \code{git} may first ask to
install the Xcode command line tools. The recursive clone and
\fname{install.sh} fetch a few gigabytes: be patient. The dot before
\fname{export.sh} matters, because it makes the script change \emph{this}
terminal. The procedure works on Apple Silicon and on Intel.

\shfile[caption={The toolchain on macOS.},label=lst:setup:mac]{code/l01/c10.sh}

\subsection{Linux (Debian, Ubuntu)}

The prerequisites are one long \code{apt} line, to be copied whole
(Listing~\ref{lst:setup:linux}); steps 2 and 3 are the same as on macOS. Other
distributions have the same packages under other names. For the kit only, the
user must belong to the group \code{dialout}, which gives access to the serial
port of the board: \code{sudo usermod -aG dialout \$USER}, then log out and in
again.

\shfile[caption={The toolchain on Linux.},label=lst:setup:linux]{code/l01/c11.sh}

\subsection{Every new terminal: switch ESP-IDF on}

\fname{export.sh} acts on the terminal in which it runs, and only on that one.
A new window is a new terminal, and must be activated again. Two aliases save
typing (Listing~\ref{lst:setup:activate}); they are a convenience, not a
requirement. The \code{prg32} commands are run from the \prg folder.

\shfile[caption={Activating ESP-IDF in a new terminal.},label=lst:setup:activate]{code/l01/c12.sh}

\begin{prgtry}[See the error once, on purpose]
Close the terminal, open a new one and run \code{python3 -m prg32 doctor}
\emph{without} activating ESP-IDF. Read the message. It is better to meet it
here than in the middle of a tutorial.
\end{prgtry}

\subsection{The doctor}

The command \code{doctor} checks the installation
(Listing~\ref{lst:setup:doctor}). The line about \code{idf.py} says that
ESP-IDF is active; the one about \code{riscv32-esp-elf-gcc} that the \riscv
compiler is reachable; those about partitions that the \prg folder is
complete. The four cartridge slots, \code{cart0} to \code{cart3}, are regions
of the flash memory of the board, of 131\,072 bytes each. Any line marked
\code{ERROR} says what to do. If the doctor passes, level~2 is done.

\txtfile[caption={The doctor (messages of October 2026).},label=lst:setup:doctor]{code/l01/c13.sh}

% ----------------------------------------------------------------------------
\section{Step 3: your first build}\label{sec:setup:build}

\subsection{Build a cartridge from the examples}

The repository contains complete games. Building one of them proves that the
compiler works before a single line has been typed
(Listing~\ref{lst:setup:build}).

\txtfile[caption={Building Pong, in C, from the examples.},label=lst:setup:build]{code/l01/c14.sh}

The command \code{cartridge build} takes a source and produces a
\fname{.prg32} file. The option \code{-{}-entry-prefix} gives the name that
precedes \code{\_init}, \code{\_update} and \code{\_draw} in the source;
\code{-{}-portable} asks for a cartridge that runs on every host, and not on
one firmware only. A whole Pong is three thousand bytes. In
Chapter~\ref{ch:firstcontact} the source will be the reader's own, in
assembly, and the command will be the same. On Windows write \code{python}
and put the command on one line.

\subsection{Send it to PRG32-QT}

\prgqt answers on port 8080 exactly as a board does
(Listing~\ref{lst:setup:send}). The command \code{runtime} prints the state of
the host and is a quick test of the connection; \code{upload-and-run} stores
the cartridge in a slot and starts it. The command says \code{esp32c6}
because, for the tool, \prgqt \emph{is} a board. If \code{localhost} does not
answer, read the address shown on the Setup screen of \prgqt and use it.
Without the tools, the same result is obtained in \prgqt by importing the file
\fname{pong.prg32} and choosing \emph{Run Cartridge}; with
\code{-{}-url http://192.168.4.1} the same command sends the game to a real
board.

\txtfile[caption={Sending the cartridge to \prgqt, running on the same computer.},label=lst:setup:send]{code/l01/c15.sh}

\subsection{Look inside: the debugger}

\begin{prgtry}[Open the machine]
\begin{enumerate}[nosep]
\item In \prgqt open \emph{Settings} and enable \emph{Run cartridges in RISC-V
  debug mode}.
\item Run the cartridge again: the game is on the left, the machine on the
  right (Figure~\ref{fig:setup:debug}).
\item Press \key{Pause} (the button with two bars): the game freezes, and one
  instruction is highlighted in the listing.
\item Press \key{Step} (the triangle with a bar) a few times: the highlight
  moves, and some registers change value.
\item Press \key{Update}, then \key{Draw}: the listing jumps to those two
  functions of the cartridge.
\item Press \key{Resume} (the triangle): the game goes on. Lower the speed to
  0.1x and watch it in slow motion.
\end{enumerate}
\end{prgtry}

The three transport buttons of the debugger are drawn as on a video player;
this book calls them \key{Pause}, \key{Resume} and \key{Step}.
This is only a first look: the next chapter explains what the registers are.
The point here is that the debugger opens and responds. Debug mode stays on
until it is switched off.

\begin{figure}[htb]
\centering
\includegraphics[width=0.95\textwidth]{prg32qt_debugger_details.png}
\caption{The machine, open: assembly with the current instruction, the 32
registers, memory, and the counters of the frame. If you see this with your
own cartridge, you are ready for the rest of the book.}
\label{fig:setup:debug}
\end{figure}

% ----------------------------------------------------------------------------
\section{Optional: QEMU and the kit}

\subsection{Level 3: the QEMU emulator}

QEMU emulates a whole chip and runs the real firmware
(Listing~\ref{lst:setup:qemu}). The command \code{qemu build} compiles the
firmware and takes several minutes the first time, because it downloads
components; \code{qemu upload} stages a cartridge into the flash image, and
QEMU must be closed while it runs. If the upload fails, run QEMU once so that
the flash image exists, close it, and upload again. Since October 2026 the
firmware starts on the \emph{Setup} menu by default, on QEMU and on the
board: choose \emph{Run Cartridge}, then the slot \code{cart0}. This procedure comes from
the documentation of \prg (\fname{docs/usage/qemu.md}) and was not run for
this book.

\shfile[caption={Level 3: the QEMU emulator.},label=lst:setup:qemu]{code/l01/c18.sh}

\subsection{Level 4: the kit}

The firmware is written to the board once, through a USB cable; after that,
games arrive over Wi-Fi and no cable is needed
(Listing~\ref{lst:setup:kit}). The upload command is the same as for \prgqt,
with another address. If the serial port is not found it is passed
explicitly: \code{COM}$x$ on Windows, \fname{/dev/ttyACM0} on Linux,
\fname{/dev/cu.usbmodem\ldots} on macOS. A black screen after flashing usually
means that the firmware for QEMU was flashed instead of the one for the board.
This procedure, too, comes from the documentation and was not run for this
book.

\shfile[caption={Level 4: the kit.},label=lst:setup:kit]{code/l01/c19.sh}

% ----------------------------------------------------------------------------
\section{When something does not work}

Table~\ref{tab:setup:trouble} lists the symptoms met most often; its first row
covers most cases. The file \fname{docs/usage/troubleshooting.md} of the
repository has the complete list.

\begin{table}[htb]
\centering\small
\caption{Troubleshooting.}\label{tab:setup:trouble}
\begin{tabularx}{\textwidth}{@{}P{0.27\textwidth}LL@{}}
\toprule
\textbf{You see} & \textbf{It means} & \textbf{Do this}\\
\midrule
\code{idf.py: command not found} & ESP-IDF is not active in this terminal & run \fname{export.sh}, or open the ESP-IDF PowerShell\\
\code{riscv32-esp-elf-gcc} missing & the compilers were not installed, or not activated & \code{./install.sh esp32c3,esp32c6}, then \fname{export.sh}\\
\code{No module named prg32} & you are not in the \prg folder & \code{cd} to the folder that you cloned\\
upload: connection refused & \prgqt is not running, or a firewall blocks port 8080 & start it; allow it; use the address on its Setup screen\\
macOS will not open PRG32 & the application is from outside the App Store & right-click, Open; or Privacy and Security, Open Anyway\\
Linux: PRG32 does not start & the Qt of the system is older than 6.5 & a newer distribution, or build from source\\
QEMU ignores the keys & the keyboard goes to the terminal & click the terminal that launched QEMU\\
\bottomrule
\end{tabularx}
\end{table}

\begin{prgconcept}[How to ask for help]
Paste the command and the \emph{whole} message as text. Never send a
photograph of the screen.
\end{prgconcept}

\subsection{Ready for the next chapter?}

Five checks, and the workshop is ready.

\begin{enumerate}
\item \prgqt starts and runs a game from the Store (step~1).
\item In a new terminal, after activating ESP-IDF, \code{doctor} shows only
  \code{OK} (step~2).
\item \code{cartridge build} produces \fname{pong.prg32} (step~3).
\item \code{upload-and-run} starts it in \prgqt (step~3).
\item In debug mode, \key{Pause} and \key{Step} work (step~3).
\end{enumerate}

% ----------------------------------------------------------------------------
% Generated by tools/textbook/export_listings.js from decks/plan.js and code/l01.S: do not edit.
\section{The program of this chapter}\label{sec:lab:l01}

The listings of this chapter are parts of one program, \fname{code/l01.S} in the
repository of the course, printed in full in Listing~\ref{lst:lab:l01}. The
automatic tests of the course material build it and run it on the headless
runner of \prgqt{}, with scripted input, and check the frame;
Figure~\ref{fig:lab:l01} is the frame that the test wrote.

\begin{description}[style=nextline,leftmargin=1.5em]
\item[On the screen.] Four vertical bars, red, green, blue and white, and the words PRG32 OK.
\item[What to press.] Nothing: the program runs by itself.
\end{description}

\begin{figure}[htb]
\centering
\includegraphics[width=0.5\textwidth]{lab_l01.png}
\caption{\fname{l01.S} on \prgqt{}.}\label{fig:lab:l01}
\end{figure}

\asmfile[caption={The complete program of the chapter, \fname{code/l01.S}.},label=lst:lab:l01]{code/l01/program.S}

To run it on \prgqt{} the commands are those of Listing~\ref{lst:lab:l01:run};
Section~\ref{sec:prg:testing} gives the same steps for QEMU and for the board.

\begin{shell}[caption={Building \fname{l01.S} and running it on \prgqt{}.},label=lst:lab:l01:run]
# 1. start PRG32-QT on this computer (it listens on port 8080)

# 2. in the PRG32 folder, with ESP-IDF active: build the cartridge
python3 -m prg32 cartridge build ~/arc12/code/l01.S --portable \
    --entry-prefix l01 --name l01 --out l01.prg32

# 3. send it to PRG32-QT and start it
python3 -m prg32 esp32c6 upload-and-run l01.prg32 \
    --url http://localhost:8080 --slot cart0

# 4. look, press, compare with the previous slide
\end{shell}


% ----------------------------------------------------------------------------
\section{Check your understanding}

\begin{quiz}{In a new terminal, \code{python3 -m prg32 doctor} says:
\code{idf.py missing}. What is the most likely cause?}
\item ESP-IDF must be installed again
\item ESP-IDF is not active in this terminal
\item \prgqt is not running
\item the board is not connected
\end{quiz}

\begin{quiz}{Which host lets you run and debug a cartridge without installing a compiler?}
\item the ESP32-C6 board
\item QEMU
\item \prgqt
\item none of them
\end{quiz}

\begin{quiz}{Why does the same \fname{pong.prg32} run on the board, on QEMU and on \prgqt?}
\item it is recompiled at each start
\item it contains three programs
\item it calls the runtime through the same table of functions on every host
\item \prgqt translates it to C
\end{quiz}

\section{Exercises}

\begin{exercise}[the workshop survives a restart]
\begin{enumerate}[nosep]
\item Finish the five checks above on your own laptop.
\item Download three cartridges from the Store, and find in each the first
  instruction of \code{update} with the debugger.
\item Build one of the other examples in \fname{examples/games}, and run it.
\item Close everything, restart the computer, and get back to a running
  cartridge without this book.
\item Optional: install QEMU and run the same cartridge there.
\end{enumerate}
\expected{Your workshop survives a restart, and you can open it without help.
The fourth task is the real test: it shows whether the activation of ESP-IDF
has been understood.}
\end{exercise}

\section{Glossary}

\begin{glossarytable}
toolchain & catena di strumenti di sviluppo\\
cartridge & cartuccia\\
emulator & emulatore\\
to build & compilare (costruire)\\
host & macchina ospite\\
firmware & firmware\\
runtime & ambiente di esecuzione\\
interpreter & interprete\\
to flash & scrivere nella memoria flash\\
breakpoint & punto di interruzione\\
\end{glossarytable}

\section{Summary and reading}

\prg is a runtime, a cartridge format, a tool and a Store; a cartridge runs
unchanged on the board, on QEMU and on \prgqt. \prgqt is enough to play and to
debug; ESP-IDF and the \prg repository are needed to build. ESP-IDF must be
activated in every new terminal: remember this and most problems disappear.
You have built a cartridge, sent it to \prgqt and stopped it in the debugger:
this is the loop of the whole book.

\paragraph{Reading.} The \fname{README} and the folder \fname{docs/usage} of
the \prg repository \citep{prg32repo}; the \fname{README} and the
documentation of \prgqt \citep{prg32qt}; the chapters on setup and on \prgqt
of \emph{Playing with RISC-V} \citep{prg32textbook}. Appendix~\ref{app:prg32}
of this book is a reference of the platform and of its commands.
